\section{Reduction to a dlt fibre over a curve}
\label{sec:curve}

Our objective is the integrality condition~\eqref{eq:relative-coefficient-target}.
Fix a prime $P\subset W$.  Take a smooth projective model $V\to X$ mapping
to $W$ by $g$, with SNC markings for the crepant boundary $B_V^c$, the
exceptional divisors over $X$, and the support of $g^*P$.
The morphism $g$ has connected fibres: its function field extension is
$\C(W)=\C(Z)\subset\C(V)=\C(X)$, which is regular. Its Stein
factor is therefore finite birational over the normal variety $W$ and
hence is $W$ itself.  Let $\theta_V$ be the rational top form compatible
with the chosen canonical divisor on $X$, and put $K_V=\Div(\theta_V)$.
Crepant pullback gives
\begin{equation}\label{eq:relative-resolution-presentation}
 K_V+B_V^c+\frac1{p_0}\Div(\psi)=g^*D_W.
\end{equation}
Put $t=t_P$.  At the generic point of $P$, the boundary
$B_V^c+t g^*P$ has coefficients at most one, with equality on at least one
component dominating $P$.  This is the SNC threshold calculation.

Choose sufficiently general very ample divisors $A_1,\ldots,A_{k-1}$ on
$W$ and put $C=A_1\cap\cdots\cap A_{k-1}$, taking $C=W$ when $k=1$.
Bertini, applied successively to the pullback linear systems and the
finitely many marked strata, makes $C$ and $V_C=g^{-1}C$ smooth, with
SNC restricted markings.  The curve $C$ is connected, and the restriction
of $g$ has connected fibres, so $V_C$ is connected and hence integral.
We also require $C\not\subset\Supp D_W$ and
$V_C\not\subset\Supp K_V\cup\Supp B_V^c\cup\Supp\Div(\psi)$.
The dimension of $V_C$ is $d-k+1\leq4$.  Choose an intersection point
$c\in C\cap P$ transverse and general on $P$.  We may avoid at $c$ all
other components of $\Supp D_W$, the bad strata loci on $P$, and the
images of all relevant nondominating strata which do not contain $P$.

Let $h:V_C\to C$, $F=h^*[c]$, and $B^c=B_V^c|_{V_C}$.  We specify the
canonical divisor in complete-intersection adjunction in order to retain
the coefficient of $P$.  Choose general representatives $A_j'\sim A_j$
avoiding $c$, and rational functions $\chi_j\in\C(W)^*$ with
$\Div(\chi_j)=A_j'-A_j$.  Define $\theta_C$ as the iterated Poincar\'e
residue, in the order of the cuts, of
\[
 \theta_V\prod_{j=1}^{k-1}g^*\chi_j
\]
along $g^{-1}A_1,\ldots,g^{-1}A_{k-1}$, taking each residue on the
preceding cut.  The generality of the divisors makes this a nonzero
rational top form on $V_C$.  For $k=1$ the product is empty and the residue
is the identity.  Each $\chi_j$ supplies a simple pole along the cutting
divisor and a zero along its chosen representative.  Cancelling the cutting
divisors before restriction, the residue formula gives an equality of
actual divisors:
\[
\begin{aligned}
 K_{V_C}:=\Div(\theta_C)
     &=\left(\Div\left(\theta_V\prod_jg^*\chi_j\right)
                   +\sum_jg^*A_j\right)|_{V_C}\\
     &=\left(K_V+\sum_jg^*A_j'\right)|_{V_C}.
\end{aligned}
\]
Set $\psi_C=\psi|_{V_C}$ and $D_C'=(D_W+\sum_jA_j')|_C$.  All these
restrictions are defined by generality.  Restricting
\eqref{eq:relative-resolution-presentation} now gives
\begin{equation}\label{eq:relative-curve-presentation}
 K_{V_C}+B^c+\frac1{p_0}\Div(\psi_C)=h^*D'_C,
 \qquad \operatorname{coeff}_c D'_C=\alpha.
\end{equation}
The coefficient assertion follows because $C$ meets $P$ transversely at
$c$, while the $A_j'$ and all other components of $\Supp D_W$ avoid $c$.

Along $F$, the coefficients of $B^c+tF$ are at most one, with equality on a
component of $F$.  Indeed, the marked divisors dominating $P$ restrict
transversely and reducedly, and near the selected point
$g^*P|_{V_C}=F$.  Over the generic point of $C$ there is still a projective
birational comparison with the normal fibre of the original model $X$,
and the restricted exceptional divisors are exceptional there.  To justify
this last assertion, the generic point of $C$ lies where $W\to Z$ is an
isomorphism.  Generic flatness and openness of geometric normality and
integrality give normal integral original fibres on an open base, and the
birational isomorphism locus is fibrewise dense there.  Shrinking further
makes the exceptional images have fibre codimension at least two.

On $V_C$ define an effective boundary $\Gamma$ by taking $F_{\mathrm{red}}$,
the horizontal strict transforms of the original boundary with their
original coefficients, and all horizontal restricted exceptional divisors
with coefficient one.  Then $(V_C,\Gamma)$ is log smooth, its horizontal
coefficients belong to the fixed set $\Phi\cup\{1\}$, and
\begin{equation}\label{eq:relative-error}
 E=\Gamma-B^c-tF
\end{equation}
is effective near $c$ and has coefficient zero on a component of $F$.
On the generic fibre $E$ is effective: its horizontal nonexceptional
terms cancel, while every horizontal exceptional term has coefficient
$1-\operatorname{coeff}B^c\geq0$.  Its support there is exceptional over
the original normal fibre by the preceding birational comparison.
Equation~\eqref{eq:relative-curve-presentation} gives the exact identity
\[
 K_{V_C}+\Gamma+\frac1{p_0}\Div(\psi_C)
     =h^*(D'_C+t[c])+E.
\]

We run a minimal-model program over $C$.  All negative coefficients of
$E$ are vertical and occur over finitely many points, so there is an
effective rational divisor $A$ on $C$ with $E^+:=E+h^*A\geq0$.  The
preceding identity becomes
\[
 K_{V_C}+\Gamma+\frac1{p_0}\Div(\psi_C)
     =h^*(D'_C+t[c]-A)+E^+.
\]
In particular, $K_{V_C}+\Gamma\sim_{\Q,C}E^+$; its restriction to a very
general fibre is rationally linearly equivalent to an effective divisor,
so it is pseudo-effective over $C$.

The existence and preservation results recalled in
\cite[Remark~2.11, arXiv version~2]{ChenTsakanikas2023} permit an MMP with
scaling of an ample divisor over the base for an lc generalized pair whose
underlying variety is $\Q$-factorial klt.  We apply this to the
$\Q$-factorial dlt pair $(V_C,\Gamma)$ over $C$; the same reference states
that its stages remain $\Q$-factorial and dlt.  Each birational step is
negative for the transformed adjoint.  The zero nef datum satisfies the
source's NQC hypothesis.  In dimension four,
\cite[Theorem~1.1, arXiv version~2]{ChenTsakanikas2023} terminates every
sequence of flips over $C$ for an NQC lc generalized pair whose adjoint is
pseudo-effective over $C$.  The three-dimensional statement
\cite[Theorem~1.2, arXiv version~2]{ChenTsakanikas2023} has no
pseudo-effectivity hypothesis.

These results supply the program and terminate a possible flip tail.
We check that its endpoint is nef and retains the exact identity.  Use
$\theta_C$ to choose canonical divisors on every birational stage
$h_i:V_i\to C$, and let $\Gamma_i$ and $E_i^+$ denote the pushforwards.
Every step is over $C$, so the preceding identity pushes forward to
\[
 K_{V_i}+\Gamma_i+\frac1{p_0}\Div(\psi_C)
     =h_i^*(D'_C+t[c]-A)+E_i^+,\qquad E_i^+\geq0.
\]
Thus the adjoint remains pseudo-effective over $C$.  A fibre-type negative
extremal contraction is impossible: a general positive-dimensional
contracted fibre has a curve $\ell$ not contained in $\Supp E_i^+$.
Since $V_i$ is $\Q$-factorial, a Cartier multiple of $E_i^+$ restricts
to an effective divisor on the normalization of $\ell$.  Hence
\[
 0\leq E_i^+\cdot\ell=(K_{V_i}+\Gamma_i)\cdot\ell<0,
\]
a contradiction.  Divisorial contractions strictly decrease the relative
Picard number, so an infinite program would eventually consist only of
flips.  The cited termination statements exclude such a tail in dimensions
three and four; in dimensions at most two no flips occur.  The endpoint is
a projective morphism $f_N:N\to C$ with $N$ $\Q$-factorial,
$(N,\Gamma_N)$ dlt, and $K_N+\Gamma_N$ nef over $C$.  Its unchanged regular
function-field extension and Stein factorization give connected fibres,
so $f_N$ is a contraction.

Let $E_N^+$ be the transform of $E^+$ on $N$.  The pushforward of
$h^*A$ is $f_N^*A$, so
$E_N:=E_N^+-f_N^*A$ is the pushforward of the original $E$.  Substituting
$E_N^+=E_N+f_N^*A$ into the transformed identity gives
\begin{equation}\label{eq:relative-output-error}
 K_N+\Gamma_N+\frac1{p_0}\Div(\psi_C)
   =f_N^*(D'_C+t[c])+E_N.
\end{equation}

We claim that $E_N=0$ near $c$.  Take a common resolution with maps
$r:\widetilde V\to V_C$ and $r':\widetilde V\to N$.  Discrepancy
comparison for the $K+\Gamma$ MMP gives
\[
 r^*(K_{V_C}+\Gamma)-r'^*(K_N+\Gamma_N)\geq0
\]
by \cite[Lemmas~3.38--3.39]{KollarMori1998}.  Pulling back the exact input
and output identities cancels the common base pullbacks and the common
principal divisor, giving
\begin{equation}\label{eq:relative-error-monotonicity}
 r'^*E_N\leq r^*E.
\end{equation}
Let $\eta$ be the generic point of $C$ and
$q:\widetilde V_\eta\to X_\eta$ the morphism to the normal original fibre.
The restriction of $E_N$ to $N_\eta$ is effective, being the pushforward
of the effective restriction of $E$; hence the restriction of $r'^*E_N$
is effective as well.  The restriction of $r^*E$ is effective and
$q$-exceptional, so \eqref{eq:relative-error-monotonicity} makes the
restriction of $r'^*E_N$ $q$-exceptional.  Finally,
\eqref{eq:relative-output-error} gives
$E_N\equiv_C K_N+\Gamma_N$, so that restriction is nef, in particular
$q$-nef.  The negativity lemma makes it zero.  Thus $E_N$ has no
horizontal component.  After shrinking about $c$, it is effective and
supported on the special fibre, and it remains nef over $C$.

Write
\[
 F_N=f_N^*[c]=\sum_i e_iS_i,\qquad E_N=\sum_i q_iS_i,
 \qquad e_i>0,\quad q_i\geq0,
\]
and set $\lambda=\max_i(q_i/e_i)$.  The divisor
$\Delta=\lambda F_N-E_N$ is effective, anti-nef over $C$, and misses at
least one component of the fibre.  Suppose the fibre has positive
dimension and $\Delta\ne0$.  Connectedness gives a component $S$ outside
$\Supp\Delta$ that meets it.  Since $N$ is $\Q$-factorial, choose a
positive integer $b$ so that $b\Delta$ is Cartier.  Its restriction to the
integral projective component $S$ is a nonzero effective Cartier divisor.
If $n=\dim S$ and $A_S$ is a very ample
Cartier divisor on $S$, then
\[
 \bigl((b\Delta)|_S\bigr)\cdot A_S^{n-1}>0.
\]
Here $A_S^{n-1}$ is represented by an effective curve cycle in the fibre
(by $S$ itself when $n=1$), so anti-nefness gives the opposite inequality.
This contradiction proves $\Delta=0$, hence $E_N=\lambda F_N$.
For a zero-dimensional fibre, $f_N$ is an isomorphism of normal curves
and the same proportionality is immediate.  Finally,
\[
 r'^*F_N=r^*F,
\]
and the original $E$ has coefficient zero on a component of $F$.
Taking the coefficient of its strict transform in
\eqref{eq:relative-error-monotonicity} forces $\lambda=0$, whether or not
that component survives on $N$.
This proves the claim.

Shrink $C$ about $c$ and set $T=(f_N^*[c])_{\mathrm{red}}$.  No divisor was
extracted in the program, so every surviving component of this fibre has
coefficient one in $\Gamma_N$.  Thus
$\Gamma_N=T+H$, where $H$ is horizontal and has coefficients in the fixed
set $\Phi\cup\{1\}$.  The reduced fibre $T$ is connected and projective.
Its components have dimension $d-k\leq3$, since they are the components
of the effective Cartier divisor $f_N^*[c]$ on the integral variety $N$.
Removing all boundary from the $\Q$-factorial
dlt pair shows that $N$ is klt.  Shrink again so that $D'_C$ is supported
at most at $c$.  For the integral Weil divisor
$L=p_0(K_N+T+H)$, the exact relation is now
\begin{equation}\label{eq:relative-local-cyclic}
 L+\Div(\psi_C)=p_0(\alpha+t)f_N^*[c].
\end{equation}
